\chapter{Oprogramowanie bez systemu operacyjnego (Bare-metal)}
\label{ch:bare_metal}

Na wyizolowanym uprzednio rdzeniu procesora (Cortex-A55) wykonuje się krytyczny czasowo kod aplikacji w~trybie bare-metal. Działa on całkowicie niezależnie, posiadając bezpośredni dostęp do zasobów sprzętowych. Rozwiązanie to całkowicie eliminuje narzut wydajnościowy (ang. \textit{overhead}) wynikający z~zarządzania wątkami, planowania procesów (ang. \textit{scheduling}) czy obsługi pamięci wirtualnej przez system operacyjny Linux.

Pełny kod źródłowy aplikacji bare-metal opracowanej na potrzeby niniejszej pracy dostępny jest w~publicznym repozytorium pod adresem: \url{https://github.com/GM-Benji/agilex-bare-metal.git}.

\section{Punkt wejścia i~inicjalizacja}
Uruchomienie rdzenia jest inicjowane przez sprzętowy sygnał resetu, który wymusza skok do wektora startowego (ang. reset vector). Pierwszym etapem jest wykonanie niskopoziomowego kodu startowego (napisanego w~języku asembler), który odpowiada za fundamentalną konfigurację sprzętową: poprawne ustawienie wskaźnika stosu (ang. \textit{Stack Pointer}) oraz inicjalizację tablicy wektorów przerwań. Następnie sterowanie jest przekazywane do głównej funkcji programu (\texttt{main}), zaimplementowanej w~języku C.

Zanim aplikacja przejdzie do nasłuchiwania poleceń, realizowana jest faza inicjalizacji pamięci współdzielonej (IPC). Kod bare-metal proaktywnie zeruje cały blok pamięci DRAM przeznaczony na bufory kołowe, co zapobiega interpretacji przypadkowych, resztkowych danych w~pamięci jako błędnych poleceń od systemu Linux. Następnie resetowane są wewnętrzne wskaźniki buforów (\texttt{head} oraz \texttt{tail}).

\section{Pętla główna i~wymiana danych IPC}
Wymiana informacji z~głównym systemem operacyjnym opiera się na koncepcji asynchronicznego odpytywania (ang. \textit{polling}). Główna pętla aplikacji cyklicznie sprawdza stan niezależnych buforów kołowych, co stanowi bezpośrednie odzwierciedlenie mechanizmów opisanych w~rozdziale dotyczącym architektury systemu.

Aby zrealizować bezkolizyjny odczyt, rdzeń bare-metal monitoruje bufor \texttt{linux\_to\_bm}. Jeśli wartość flagi nadawczej \texttt{head} (kontrolowanej przez Linux) jest różna od wartości flagi odbiorczej \texttt{tail} (kontrolowanej przez bare-metal), oznacza to pojawienie się nowych danych. Po odczytaniu paczki danych, kod przesuwa wskaźnik \texttt{tail}, informując tym samym nadawcę o~zwolnieniu miejsca w~buforze. W~odpowiedzi aplikacja formułuje i~odsyła komunikat zwrotny (ang. \textit{echo}) poprzez bliźniaczy bufor \texttt{bm\_to\_linux}.

\begin{lstlisting}[language=C, caption={Główna pętla asynchroniczna aplikacji bare-metal}, label={lst:bm_main}]
while (1) {
    if (shared_ram->linux_to_bm.head != shared_ram->linux_to_bm.tail) {
        int bytes = read_from_ring(&shared_ram->linux_to_bm, rx_buf, 255);
        if (bytes > 0) {
            write_to_ring(&shared_ram->bm_to_linux, "ACK: ", 5);
            write_to_ring(&shared_ram->bm_to_linux, rx_buf, bytes);
        }
    }
}
\end{lstlisting}

\section{Bariery pamięci w~architekturze wielordzeniowej}
Niezwykle istotnym aspektem implementacji komunikacji \textit{lock-free} (bez użycia muteksów czy semaforów) w~systemach opartych na architekturze ARM jest zarządzanie spójnością pamięci podręcznej. Z~uwagi na to, że rdzenie mogą w~sposób niezależny optymalizować i~zmieniać kolejność wykonywania instrukcji zapisu, wprowadzono sprzętowe bariery pamięci (ang. \textit{Data Memory Barrier}).

W~napisanych funkcjach obsługujących zapis i~odczyt (\texttt{write\_to\_ring} oraz \texttt{read\_from\_ring}) wykorzystano asemblerową instrukcję \texttt{dmb sy}. Gwarantuje ona, że bajt danych fizycznie dotrze do współdzielonej pamięci DRAM \textit{zanim} zaktualizowana zostanie wartość flagi \texttt{head} lub \texttt{tail}. Zapobiega to krytycznym błędom współbieżności, w~których odbiorca mógłby zareagować na przesunięcie wskaźnika, zanim faktyczny komunikat zostałby fizycznie utrwalony na magistrali pamięci.

\begin{lstlisting}[language=C, caption={Wykorzystanie barier pamięci (DMB) podczas zapisu do bufora kołowego}, label={lst:bm_ring}]
int write_to_ring(RingBuffer *rb, const uint8_t *buf, uint32_t len) {
    uint32_t i, next;
    for (i = 0; i~< len; i++) {
        next = (rb->head + 1) % RING_BUFFER_SIZE;
        if (next == rb->tail) break;
        rb->data[rb->head] = buf[i];
        __asm__ volatile ("dmb sy" ::: "memory");
        rb->head = next;
        __asm__ volatile ("dmb sy" ::: "memory");
    }
    return i; 
}
\end{lstlisting}